% Source: https://artofproblemsolving.com/community/c3h2492319_100th_post_problems
% Stable identity: aops-topic-2492319-q05
With probability $1$ the four points are distinct, no two are diametrically opposite, and they form a convex quadrilateral $Q$ inscribed in the circle with centre $O$. We may ignore these probability-$0$ exceptions.

\textbf{When does $Q$ contain $O$?} If all four points lie in an open semicircle, the diameter bounding it has the whole of $Q$ strictly on one side, so $O\notin Q$. Conversely, if $O\notin Q$, then since $Q$ is convex there is a line through $O$ with $Q$ strictly on one side of it; the four points then lie in the open semicircle on that side. So $Q$ contains $O$ exactly when the points do \emph{not} all lie in a semicircle.

\begin{asy}
usepackage("amsmath");
size(6cm);
defaultpen(fontsize(10pt));
void pic(pair c, real[] a, bool inside){
  draw(circle(c,1),linewidth(0.8)); dot(c); label("$O$",c,dir(-60),fontsize(8pt));
  pair[] p; for(int i=0;i<4;++i) p.push(c+dir(a[i]));
  fill(p[0]--p[1]--p[2]--p[3]--cycle,inside?rgb(0.8,0.9,1):rgb(1,0.88,0.85));
  draw(p[0]--p[1]--p[2]--p[3]--cycle,inside?blue:red);
  for(int i=0;i<4;++i) dot(p[i]);
  dot(c);
}
pic((0,0),new real[]{20,110,200,290},true);
pic((2.8,0),new real[]{10,60,120,165},false);
draw((2.8,0)+1.15*dir(0)--(2.8,0)+1.15*dir(180),dashed);
label("contains $O$",(0,-1.2),S,fontsize(8pt));
label("in a semicircle",(2.8,-1.2),S,fontsize(8pt));
\end{asy}

\textbf{Probability that all four lie in a semicircle.} For each point $P_i$ let $A_i$ be the event that the other three points lie in the semicircle that starts at $P_i$ and runs clockwise. Each of the other points lies there with probability $\dfrac12$, independently, so $P(A_i)=\dfrac18$. If the four points lie in some semicircle, then the first of them in clockwise order around that semicircle is a point $P_i$ for which $A_i$ holds; and two of the events cannot happen together (except with probability $0$), because if $A_i$ and $A_j$ both held, $P_j$ would be in the clockwise semicircle from $P_i$ and $P_i$ in the clockwise semicircle from $P_j$, forcing $P_i$ and $P_j$ to be opposite. Hence
\[P(\text{all in a semicircle})=4\cdot\dfrac18=\dfrac12,\]
and the probability that the quadrilateral contains the centre is $1-\dfrac12=\boxed{\dfrac12}$.
